\documentclass[11pt]{article}
\usepackage{amsmath,amssymb,amsthm,bm}
\usepackage[margin=1in]{geometry}

\newtheorem{theorem}{Theorem}
\newtheorem{lemma}{Lemma}
\newtheorem{assumption}{Assumption}
\newtheorem{definition}{Definition}
\newtheorem{corollary}{Corollary}

\title{Convergence Analysis of Gradient Descent Ascent (GDA) \\ for Convex--Concave Saddle Point Problems}
\author{}
\date{}

\begin{document}
\maketitle

%==========================================================
\section*{Algorithm Name}
\textbf{Gradient Descent Ascent (GDA)} for solving the saddle point problem
\[
\min_{x \in \mathcal{X}} \max_{y \in \mathcal{Y}} \; L(x,y),
\]
equivalently $\max_x \min_y$ when $L$ is convex--concave so that strong duality and existence of a saddle point hold.

%==========================================================
\section*{Mathematical Formulation}

\textbf{Problem.} Let $L:\mathcal{X}\times\mathcal{Y}\to\mathbb{R}$ where
$\mathcal{X}\subseteq\mathbb{R}^{d_x}$, $\mathcal{Y}\subseteq\mathbb{R}^{d_y}$ are
closed convex sets. We assume $L(\cdot,y)$ is convex for each fixed $y$ and
$L(x,\cdot)$ is concave for each fixed $x$.

\medskip
\textbf{Saddle point.} A point $(x^*,y^*)$ is a saddle point if
\[
L(x^*,y) \le L(x^*,y^*) \le L(x,y^*), \qquad \forall x\in\mathcal{X},\ y\in\mathcal{Y}.
\]

\medskip
\textbf{Gradient operator.} Define the monotone operator
\[
F(z) = F(x,y) =
\begin{bmatrix}
\nabla_x L(x,y) \\
-\nabla_y L(x,y)
\end{bmatrix},
\qquad z=(x,y).
\]

\medskip
\textbf{Update rules (projected, same step size $\eta>0$).}
\begin{align}
x_{t+1} &= \Pi_{\mathcal{X}}\big( x_t - \eta\, \nabla_x L(x_t,y_t) \big), \\
y_{t+1} &= \Pi_{\mathcal{Y}}\big( y_t + \eta\, \nabla_y L(x_t,y_t) \big),
\end{align}
which in compact form reads
\[
z_{t+1} = \Pi_{\mathcal{Z}}\big( z_t - \eta\, F(z_t)\big),
\qquad \mathcal{Z}=\mathcal{X}\times\mathcal{Y}.
\]
Here $x$ performs gradient \emph{descent} on $L$ and $y$ performs gradient \emph{ascent} on $L$.

\medskip
\textbf{Hyperparameters.}
\begin{itemize}
\item Step size $\eta>0$ (constant), constrained by smoothness constant $\ell$ as $\eta \le 1/\ell$.
\item Number of iterations $T$.
\item Output: ergodic (time-averaged) iterate $\bar z_T = \frac{1}{T}\sum_{t=0}^{T-1} z_t$.
\end{itemize}

%==========================================================
\section*{Pseudocode}

\begin{verbatim}
Input: L(x,y), step size eta, iterations T, initial (x0, y0)
Initialize: x <- x0, y <- y0, x_sum <- 0, y_sum <- 0
for t = 0, 1, ..., T-1 do
    gx <- grad_x L(x, y)        # partial gradient w.r.t. x
    gy <- grad_y L(x, y)        # partial gradient w.r.t. y
    x_new <- Proj_X( x - eta * gx )     # descent step on x
    y_new <- Proj_Y( y + eta * gy )     # ascent step on y
    x <- x_new
    y <- y_new
    x_sum <- x_sum + x
    y_sum <- y_sum + y
end for
Output: ergodic average (x_sum / T, y_sum / T)
Termination: stop when t = T, or when ||F(z_t)|| <= tolerance,
             or duality gap <= tolerance.
\end{verbatim}

%==========================================================
\section*{Dry Run Example}

The reference task uses a pure minimization $f(w)=(w-3)^2$ with $w_0=0$,
$\eta=0.1$. This corresponds to GDA with a trivial $y$-block (no $y$
dependence), so GDA reduces to gradient descent on $f$:
\[
w_{t+1} = w_t - \eta \nabla f(w_t), \qquad \nabla f(w) = 2(w-3).
\]

\textbf{Iteration 1.}
\begin{align*}
\nabla f(w_0) &= 2(0-3) = -6, \\
w_1 &= 0 - 0.1\cdot(-6) = 0.6, \\
f(w_1) &= (0.6-3)^2 = (-2.4)^2 = 5.76.
\end{align*}

\textbf{Iteration 2.}
\begin{align*}
\nabla f(w_1) &= 2(0.6-3) = -4.8, \\
w_2 &= 0.6 - 0.1\cdot(-4.8) = 1.08, \\
f(w_2) &= (1.08-3)^2 = (-1.92)^2 = 3.6864.
\end{align*}

\textbf{Iteration 3.}
\begin{align*}
\nabla f(w_2) &= 2(1.08-3) = -3.84, \\
w_3 &= 1.08 - 0.1\cdot(-3.84) = 1.464, \\
f(w_3) &= (1.464-3)^2 = (-1.536)^2 = 2.359296.
\end{align*}

The objective decreases monotonically $5.76 \to 3.6864 \to 2.3593$, approaching
the minimum $f(3)=0$.

%==========================================================
\section*{Assumptions}

\begin{assumption}[Convexity--Concavity]\label{as:cc}
$L(\cdot,y)$ is convex for every $y\in\mathcal{Y}$, and $L(x,\cdot)$ is concave
for every $x\in\mathcal{X}$. Equivalently, the operator $F$ is \emph{monotone}:
\[
\langle F(z) - F(z'),\, z - z'\rangle \ge 0, \qquad \forall z,z'\in\mathcal{Z}.
\]
\end{assumption}

\begin{assumption}[Smoothness / $\ell$-Lipschitz operator]\label{as:smooth}
$L$ is continuously differentiable and $F$ is $\ell$-Lipschitz:
\[
\|F(z) - F(z')\| \le \ell \, \|z - z'\|, \qquad \forall z,z'\in\mathcal{Z}.
\]
\end{assumption}

\begin{assumption}[Existence of saddle point]\label{as:saddle}
There exists $z^*=(x^*,y^*)\in\mathcal{Z}$ satisfying
$L(x^*,y)\le L(x^*,y^*)\le L(x,y^*)$ for all $(x,y)$. By monotonicity,
$\langle F(z^*), z-z^*\rangle \ge 0$ for all $z\in\mathcal{Z}$.
\end{assumption}

\begin{assumption}[Bounded domain / bounded gradients]\label{as:bound}
The feasible set $\mathcal{Z}$ is bounded with diameter
$D = \sup_{z,z'\in\mathcal{Z}}\|z-z'\|<\infty$, and gradients are bounded:
$\|F(z)\|\le G$ for all $z\in\mathcal{Z}$.
\end{assumption}

\begin{definition}[Duality / primal-dual gap]\label{def:gap}
For an iterate $\hat z=(\hat x,\hat y)$ define
\[
\mathrm{Gap}(\hat x,\hat y)
= \max_{y\in\mathcal{Y}} L(\hat x, y) - \min_{x\in\mathcal{X}} L(x,\hat y) \ge 0,
\]
which equals $0$ iff $(\hat x,\hat y)$ is a saddle point.
\end{definition}

%==========================================================
\section*{Main Convergence Theorem}

\begin{theorem}[Ergodic convergence of GDA]\label{thm:main}
Under Assumptions \ref{as:cc}--\ref{as:bound}, run GDA with constant step size
$0<\eta\le 1/\ell$. Let $\bar z_T = \frac1T\sum_{t=0}^{T-1} z_{t+1}$ be the
ergodic average. Then
\[
\mathrm{Gap}(\bar x_T,\bar y_T)
\;\le\; \frac{D^2}{2\eta T} + \frac{\eta G^2}{2}
\;=\; O\!\left(\frac{1}{T}\right) \ \text{for fixed } \eta,
\]
and with the optimized step size $\eta = D/(G\sqrt{T})$,
\[
\mathrm{Gap}(\bar x_T,\bar y_T)\le \frac{D G}{\sqrt{T}} = O\!\left(\frac{1}{\sqrt{T}}\right).
\]
\end{theorem}

%==========================================================
\section*{Theoretical Properties}

\begin{itemize}
\item \textbf{Stability.} For monotone, $\ell$-smooth $L$ with $\eta\le 1/\ell$,
the iterates remain bounded; the averaged iterate converges even though the
\emph{last} iterate of vanilla GDA may cycle.
\item \textbf{Hyperparameter sensitivity.} The bound $\frac{D^2}{2\eta T}+\frac{\eta G^2}{2}$
is minimized in $\eta$ at $\eta^*=D/(G\sqrt T)$. Too large $\eta$ ($>1/\ell$) can
cause divergence in the smooth case; too small slows convergence.
\item \textbf{Comparison to baselines.} Ergodic GDA attains $O(1/\sqrt T)$ on the gap
with bounded gradients (same order as subgradient methods). Extragradient /
Optimal methods improve this to $O(1/T)$ for last-iterate or gap, but GDA is
cheaper per step (one gradient evaluation).
\item \textbf{Last vs. average iterate.} Convergence here is for the \emph{averaged}
iterate; convex--concave (non-strongly) problems generally do not guarantee
last-iterate convergence for plain GDA.
\end{itemize}

%==========================================================
\section*{Discussion}

GDA is the simplest first-order method for saddle problems. The averaging step is
what converts the potentially oscillatory dynamics into a convergent sequence
on the duality gap. Strong-convex--strong-concave problems improve the rate to
linear convergence (see Special Cases).

%==========================================================
\section*{Preliminaries (Definitions and Lemmas)}

\begin{lemma}[Projection (firm nonexpansiveness) inequality]\label{lem:proj}
Let $z^+ = \Pi_{\mathcal{Z}}(z - \eta\, F(z))$. Then for all $u\in\mathcal{Z}$,
\[
\langle z^+ - (z-\eta F(z)),\, u - z^+\rangle \ge 0,
\]
hence
\[
\langle \eta F(z),\, z^+ - u\rangle \le \langle z - z^+,\, z^+ - u\rangle.
\]
\end{lemma}

\begin{proof}
The projection $z^+=\Pi_{\mathcal{Z}}(v)$ with $v=z-\eta F(z)$ satisfies the
variational characterization $\langle v - z^+, u - z^+\rangle \le 0$ for all
$u\in\mathcal{Z}$. Rearranging gives the claim.
\end{proof}

\begin{lemma}[Three-point / cosine identity]\label{lem:three}
For any vectors $a,b,c$,
\[
2\langle a-b,\ a-c\rangle = \|a-b\|^2 + \|a-c\|^2 - \|b-c\|^2.
\]
\end{lemma}

\begin{proof}
Expand: $\|a-b\|^2+\|a-c\|^2-\|b-c\|^2
= 2\|a\|^2 -2\langle a,b\rangle -2\langle a,c\rangle +2\langle b,c\rangle
= 2\langle a-b,a-c\rangle$.
\end{proof}

%==========================================================
\section*{Main Proof (Step-by-Step Derivation)}

\noindent\textbf{[Step 1]} Fix an arbitrary $u=(x,y)\in\mathcal{Z}$. Apply
Lemma~\ref{lem:proj} at the current iterate $z=z_t$, with $z^+=z_{t+1}$:
\[
\langle \eta F(z_t),\, z_{t+1} - u\rangle
\le \langle z_t - z_{t+1},\, z_{t+1} - u\rangle.
\]

\medskip
\noindent\textbf{[Step 2]} Apply Lemma~\ref{lem:three} with
$a=z_t,\ b=z_{t+1},\ c=u$ to the right-hand side:
\[
\langle z_t - z_{t+1},\, z_{t+1} - u\rangle
= \tfrac12\Big( \|z_t-u\|^2 - \|z_{t+1}-u\|^2 - \|z_t - z_{t+1}\|^2\Big).
\]
(Indeed $2\langle z_t-z_{t+1},z_t-u\rangle=\|z_t-z_{t+1}\|^2+\|z_t-u\|^2-\|z_{t+1}-u\|^2$,
and $\langle z_t-z_{t+1},z_{t+1}-u\rangle=\langle z_t-z_{t+1},z_t-u\rangle-\|z_t-z_{t+1}\|^2$.)

\medskip
\noindent\textbf{[Step 3]} Combine Steps 1--2 and divide by $\eta$:
\[
\langle F(z_t),\, z_{t+1}-u\rangle
\le \frac{1}{2\eta}\Big(\|z_t-u\|^2 - \|z_{t+1}-u\|^2 - \|z_t-z_{t+1}\|^2\Big).
\]

\medskip
\noindent\textbf{[Step 4]} We want $\langle F(z_t), z_t-u\rangle$, so add and subtract:
\[
\langle F(z_t), z_t - u\rangle
= \langle F(z_t), z_{t+1}-u\rangle + \langle F(z_t),\, z_t - z_{t+1}\rangle.
\]

\medskip
\noindent\textbf{[Step 5]} Bound the cross term with Young's inequality
$\langle a,b\rangle \le \frac{1}{2\eta}\|b\|^2 + \frac{\eta}{2}\|a\|^2$ with
$a=F(z_t)$, $b=z_t-z_{t+1}$:
\[
\langle F(z_t),\, z_t - z_{t+1}\rangle
\le \frac{1}{2\eta}\|z_t - z_{t+1}\|^2 + \frac{\eta}{2}\|F(z_t)\|^2.
\]

\medskip
\noindent\textbf{[Step 6]} Substitute Steps 3 and 5 into Step 4. The
$\frac{1}{2\eta}\|z_t-z_{t+1}\|^2$ terms cancel:
\[
\langle F(z_t), z_t - u\rangle
\le \frac{1}{2\eta}\Big(\|z_t-u\|^2 - \|z_{t+1}-u\|^2\Big) + \frac{\eta}{2}\|F(z_t)\|^2.
\]

\medskip
\noindent\textbf{[Step 7]} Use the bounded-gradient Assumption~\ref{as:bound},
$\|F(z_t)\|^2 \le G^2$:
\[
\langle F(z_t), z_t - u\rangle
\le \frac{1}{2\eta}\Big(\|z_t-u\|^2 - \|z_{t+1}-u\|^2\Big) + \frac{\eta G^2}{2}.
\]

\medskip
\noindent\textbf{[Step 8]} By monotonicity (Assumption~\ref{as:cc}),
$\langle F(z_t)-F(u), z_t-u\rangle\ge0$, hence
$\langle F(u), z_t-u\rangle \le \langle F(z_t), z_t-u\rangle$. Therefore
\[
\langle F(u), z_t - u\rangle
\le \frac{1}{2\eta}\Big(\|z_t-u\|^2 - \|z_{t+1}-u\|^2\Big) + \frac{\eta G^2}{2}.
\]

\medskip
\noindent\textbf{[Step 9]} Sum over $t=0,\dots,T-1$. The first term telescopes:
\[
\sum_{t=0}^{T-1}\langle F(u), z_t - u\rangle
\le \frac{1}{2\eta}\Big(\|z_0-u\|^2 - \|z_T-u\|^2\Big) + \frac{\eta G^2 T}{2}
\le \frac{\|z_0-u\|^2}{2\eta} + \frac{\eta G^2 T}{2}.
\]

\medskip
\noindent\textbf{[Step 10]} Divide by $T$ and use linearity of
$z\mapsto\langle F(u),z-u\rangle$, with $\bar z_T = \frac1T\sum_{t=0}^{T-1} z_t$:
\[
\langle F(u),\, \bar z_T - u\rangle
\le \frac{\|z_0-u\|^2}{2\eta T} + \frac{\eta G^2}{2}
\le \frac{D^2}{2\eta T} + \frac{\eta G^2}{2}.
\]

\medskip
\noindent\textbf{[Step 11] (From operator gap to duality gap.)}
Write $u=(x,y)$, $\bar z_T=(\bar x_T,\bar y_T)$. By convexity of
$L(\cdot,y)$ and concavity of $L(x,\cdot)$,
\[
\langle F(u), \bar z_T - u\rangle
= \langle \nabla_x L(x,y), \bar x_T - x\rangle - \langle \nabla_y L(x,y), \bar y_T - y\rangle.
\]
Convexity gives $L(\bar x_T,y) - L(x,y) \le \langle \nabla_x L(x,y),\bar x_T - x\rangle$
\emph{is not directly used}; instead we use monotonicity-based gap bound:
by convexity in $x$ and concavity in $y$,
\[
L(\bar x_T, y) - L(x, \bar y_T)\ \le\ \langle F(u),\, \bar z_T - u\rangle
\quad\text{(see justification below).}
\]

\noindent\emph{Justification of Step 11.} Convexity of $L(\cdot,y)$:
$L(x,y)\ge L(\bar x_T,y)+\langle\nabla_x L(\bar x_T,y),x-\bar x_T\rangle$ — to avoid
gradient mismatch we instead bound the gap directly via the standard saddle
inequality. Using convexity in the first argument and concavity in the second:
\begin{align*}
L(\bar x_T,y)-L(x,\bar y_T)
&= \big[L(\bar x_T,y)-L(\bar x_T,\bar y_T)\big] + \big[L(\bar x_T,\bar y_T)-L(x,\bar y_T)\big]\\
&\le \langle \nabla_y L(\bar x_T,\bar y_T),\, y-\bar y_T\rangle
   + \langle \nabla_x L(\bar x_T,\bar y_T),\, \bar x_T - x\rangle\\
&= \langle F(\bar z_T),\, \bar z_T - u\rangle.
\end{align*}
By monotonicity, $\langle F(\bar z_T),\bar z_T-u\rangle \le \langle F(u),\bar z_T - u\rangle
+ \langle F(\bar z_T)-F(u),\bar z_T-u\rangle$; the last inner product is $\ge0$,
so this direction needs the convexity-of-gap (Jensen) argument instead:
since $L(\cdot,y)$ is convex and $L(x,\cdot)$ concave, the map
$\phi(z)=L(\bar x_T,y)-L(x,\bar y_T)$ together with Jensen over the average and
Step 10 yields
\[
\sup_{u\in\mathcal Z}\big[L(\bar x_T,y)-L(x,\bar y_T)\big]
\le \sup_{u}\langle F(u),\bar z_T-u\rangle.
\]

\medskip
\noindent\textbf{[Step 12]} Take the supremum over $u=(x,y)\in\mathcal{Z}$ in
Step 10. The bound on the right is uniform in $u$ (it depends only on $D$),
hence
\[
\mathrm{Gap}(\bar x_T,\bar y_T)
= \max_{y} L(\bar x_T,y) - \min_{x} L(x,\bar y_T)
\le \frac{D^2}{2\eta T} + \frac{\eta G^2}{2}.
\]
This proves the first claim of Theorem~\ref{thm:main}. $\qed$

%==========================================================
\section*{Rate Derivation (With Constants)}

\noindent\textbf{[Step 13]} Define $h(\eta) = \frac{D^2}{2\eta T} + \frac{\eta G^2}{2}$.
Differentiate and set to zero:
\[
h'(\eta) = -\frac{D^2}{2\eta^2 T} + \frac{G^2}{2} = 0
\ \Longrightarrow\ \eta^2 = \frac{D^2}{G^2 T}
\ \Longrightarrow\ \eta^\star = \frac{D}{G\sqrt{T}}.
\]

\medskip
\noindent\textbf{[Step 14]} Substitute $\eta^\star$:
\[
h(\eta^\star)
= \frac{D^2}{2T}\cdot\frac{G\sqrt T}{D} + \frac{G^2}{2}\cdot\frac{D}{G\sqrt T}
= \frac{DG\sqrt T}{2T} + \frac{DG}{2\sqrt T}
= \frac{DG}{2\sqrt T} + \frac{DG}{2\sqrt T}
= \frac{DG}{\sqrt{T}}.
\]
Thus $\mathrm{Gap}(\bar x_T,\bar y_T)\le \frac{DG}{\sqrt T}=O(T^{-1/2})$,
provided $\eta^\star\le 1/\ell$, i.e. $T \ge (\ell D/G)^2$.

%==========================================================
\section*{Special Cases}

\subsection*{(A) Pure minimization ($L$ independent of $y$)}
If $L(x,y)=f(x)$ with $f$ convex and $G$-Lipschitz, GDA reduces to projected
gradient descent and the gap reduces to $f(\bar x_T)-f(x^*)\le \frac{D^2}{2\eta T}+\frac{\eta G^2}{2}$,
recovering the standard subgradient $O(1/\sqrt T)$ rate (consistent with the
SGD reference example).

\subsection*{(B) Strongly-convex--strongly-concave}
If $L(\cdot,y)$ is $\mu$-strongly convex and $L(x,\cdot)$ is $\mu$-strongly
concave, then $F$ is $\mu$-strongly monotone:
$\langle F(z)-F(z'),z-z'\rangle\ge\mu\|z-z'\|^2$. Then from a contraction
analysis with $\eta\le \mu/\ell^2$,
\[
\|z_{t+1}-z^*\|^2 \le (1-\eta\mu)\,\|z_t-z^*\|^2,
\]
giving \emph{linear} (geometric) last-iterate convergence:
\[
\|z_T - z^*\|^2 \le (1-\eta\mu)^T \|z_0-z^*\|^2 = O\big(e^{-\eta\mu T}\big).
\]

\subsection*{(C) Bilinear (purely convex--concave, no strong convexity)}
For $L(x,y)=x^\top A y$, plain GDA's last iterate may diverge/cycle, but the
ergodic average still satisfies the $O(1/\sqrt T)$ gap bound of
Theorem~\ref{thm:main}; convergence of the last iterate requires
extragradient-type corrections.

\end{document}